Himpunan ukara nggambaraké struktur kang ndadèkaké ukara-ukara
iku bebarengan bener. Struktur mau diarani \emph{modèl}.
Kosok baliné, saka sawijining struktur utawa himpunan struktur,
kita bisa njupuk himpunan ukara kang disukupi déning kabèh
struktur mau. Iki \emph{teori} struktur utawa himpunan struktur
kasebut. Saben ukara kang dadi konsèkuènsi himpunan ukara mau
wis ana ing himpunané. Himpunan iku \emph{katutup}. Sacara umum,
$\Gamma$ diarani teori yèn katutup tumrap konsèkuènsi. Kita
ngarani $\Gamma$ \emph{diaksiomatisasi} déning~$\Delta$ yèn
$\Gamma$ ngemot persis kabèh ukara kang dadi konsèkuènsi~$\Delta$.

Tuladha teori matématika yaiku teori urutan linier, teori grup,
lan teori aritmetika, kayata teori kang diaksiomatisasi déning
aksioma Peano. Teori wigati uga ana ing disiplin liya. Basis
data relasional bisa dianggep teori. Ing metafisika, teori
ngenani relasi pérangan uga bisa diaksiomatisasi.

Nalika nyusun teori, salah siji pitakonan wigati yaiku apa basané
cukup èksprèsif kanggo ngrumusaké kabèh bab kang arep dirembug.
Pitakonan liya yaiku apa teoriné cukup kuwat kanggo mbuktèkaké
asil kang dikarepaké. Kanggo \emph{ngandharaké} sawijining relasi,
kita mbutuhaké formula mawa cacah variabèl bebas kang cocog.
Ing \emph{teori himpunan}, $\in$ dadi siji-sijiné simbol relasi,
nanging $x\subseteq y$ bisa diandharaké nganggo
$\lforall[u][(u\in x\lif u\in y)]$. \emph{Teori himpunan
Zermelo--Fraenkel} $\Log{ZFC}$ cukup kuwat kanggo ngandharaké
meh saben pratelan matématika lan mbuktèkaké meh saben teorema
matématika nganggo sawatara aksioma lan ranté definisi kang saya
ruwet, kayata definisi $\subseteq$.
